\documentclass[11pt]{amsart}
\usepackage{amsmath,amssymb,amsthm,mathtools}
\usepackage[margin=1in]{geometry}
\usepackage{hyperref}
\usepackage{listings}
\lstset{basicstyle=\ttfamily\footnotesize,breaklines=true,frame=single,
  showstringspaces=false,columns=fullflexible,keepspaces=true}

\newtheorem{theorem}{Theorem}[section]
\newtheorem{proposition}[theorem]{Proposition}
\newtheorem{lemma}[theorem]{Lemma}
\newtheorem{corollary}[theorem]{Corollary}
\theoremstyle{definition}
\newtheorem{definition}[theorem]{Definition}
\theoremstyle{remark}
\newtheorem{remark}[theorem]{Remark}

\newcommand{\fa}{\mathfrak a}
\newcommand{\fp}{\mathfrak p}
\newcommand{\ft}{\mathfrak t}
\newcommand{\RH}{\mathrm{RH}}
\newcommand{\bbR}{\mathbb{R}}
\DeclareMathOperator{\rank}{rank}

\title[A Second-Moment Identity and the Localized Weil Form]{A second-moment
identity for zero exponential sums and the localized Weil form}
\author{David Alejandro Trejo Pizzo}
\thanks{Independent Researcher. \texttt{dtrejopizzo@gmail.com}}


\begin{document}
\begin{abstract}
We record the exact identity
$V(d,T)=\int_{-2d}^{2d}\bigl|\sum_\gamma e^{i\gamma\xi}\bigr|^2d\xi$ for the
ordinates $\gamma$ of the zeros of $\zeta$ in a window at height $T$, and show
how bounds on $V$ bound the corresponding windowed zero sum of the Weil
quadratic form on band-limited test functions. Assuming RH, Montgomery's
pair-correlation theorem controls the low-frequency part of a weighted version
of $V$, and, by positivity of the Fej\'er kernel, the high-frequency part is
$O(dN(T))$, so that the weighted moment satisfies $V_w\ll dN(T)$. By the
theorem of Goldston and Montgomery (under RH), the uniform high-frequency form
factor, which gives the asymptotic $V_w\sim4dN(T)$ for the weighted moment, is
equivalent to an asymptotic
for the variance of primes in short intervals. A numerical appendix compares
the two sides of the identity for $80$ zeros (they agree with quadrature to
$3\times10^{-6}$) and records exploratory computations on the prime tail. No
claim about the sign of the form, or about RH, is made.
\end{abstract}
\maketitle

\tableofcontents
\newpage

\section{Introduction}
The Weil explicit formula~\cite{Weil} realizes a quadratic functional $\ft(g)=\sum_\rho
h_g(\gamma_\rho)$ over the nontrivial zeros $\tfrac12+i\gamma_\rho$, with $h_g$
the entire extension of $|\widehat g|^2$; on-line zeros contribute nonnegatively,
off-line zeros contribute indefinite quartets, and $\RH\iff\ft\succeq0$ (Weil's
criterion~\cite{Weil}). Three properties of $\ft$ are of different strength:
\begin{align*}
&\textbf{(S)}\ \ \ft\ge-C\|g\|^2\quad(\text{semibounded, uniformly in the height});\\
&\textbf{($\kappa$)}\ \ \ft\ \text{has finitely many negative squares};\\
&\textbf{(R)}\ \ \ft\succeq0\quad(\RH).
\end{align*}
Property (R) implies the other two; we prove neither (S) nor ($\kappa$).
Companion work~\cite{RP16} reduces ($\kappa$) to a relative form bound on the
prime tail and observes numerically that such a bound with a fixed cutoff
appears to fail. This paper isolates the second moment $V(d,T)$ of the zero
exponential sum over a height window, which controls the size of the windowed
zero sum of $\ft$ (Proposition~\ref{prop:parseval}), and records what is known
about $V$: an unconditional trivial bound, a conditional bound under a
hypothesis, and, for the weighted full-range version $V_w$ of Section~3, the
low-frequency part under RH by Montgomery's theorem, a bound $V_w\ll dN(T)$
under RH, and the equivalence (under RH) of the uniform high-frequency form factor with a
short-interval prime variance asymptotic. Nothing here addresses the sign of
$\ft$, hence nothing addresses RH.

\section{The localized form and the controlling second moment}
\subsection*{Definitions}
For $g\in L^2(\bbR)$ write $\widehat g(r)=\int_\bbR g(x)e^{irx}dx$,
$\tilde g(x)=\overline{g(-x)}$, $(g\star\tilde g)(s)=\int g(x+s)\overline{g(x)}dx$,
and $\|g\|$ for the $L^2(\bbR)$ norm. A function $g$ is \emph{of type $d$} if
$g\in PW_d:=\{g\in L^2(\bbR):\operatorname{supp}g\subset[-d,d]\}$; then
$g\star\tilde g$ is supported in $[-2d,2d]$, $|(g\star\tilde g)(s)|\le\|g\|^2$,
and
\begin{equation*}
|\widehat g(r)|^2=\int_{-2d}^{2d}(g\star\tilde g)(s)\,e^{irs}\,ds\qquad(r\in\bbR).
\end{equation*}
For smooth $g\in PW_d$ the explicit formula (Weil~\cite{Weil}; in the normalization
of~\cite[\S2]{RP16}) reads $\sum_\rho h_g(\gamma_\rho)=\fa[g]-\fp[g]=:\ft[g]$, where
$h_g(r)=\widehat g(r)\overline{\widehat g(\bar r)}$,
\[
\fa[g]=\frac1{2\pi}\int_\bbR\Omega(r)|\widehat g(r)|^2dr+h_g(\tfrac i2)+h_g(-\tfrac i2),
\qquad \Omega(r)=\mathrm{Re}\,\psi(\tfrac14+\tfrac{ir}2)-\log\pi,
\]
the last two terms coming from the poles of $\zeta$ at $s=0,1$, and
\[
\fp[g]=2\sum_{n\ge2}\frac{\Lambda(n)}{\sqrt n}\,\mathrm{Re}\,(g\star\tilde g)(\log n),
\]
a finite sum ($n\le e^{2d}$). The \emph{window} at height $T$ of length $H$ is
the set of zeros $\rho$ with $T<\mathrm{Im}\,\rho\le T+H$, counted with
multiplicity; $N=N(T,H)$ is their number. For the window we put
\[
S(\xi)=\sum_{\gamma}e^{i\gamma\xi},\qquad
\ft_T[g]=\sum_{\gamma}|\widehat g(\gamma)|^2,\qquad
V(d,T)=\sum_{\gamma,\gamma'}\frac{2\sin\!\bigl(2d(\gamma-\gamma')\bigr)}{\gamma-\gamma'}
\]
(the terms with $\gamma=\gamma'$ being $4d$), all sums over $\gamma=\mathrm{Im}\,\rho$
for $\rho$ in the window. For zeros on the critical line,
$|\widehat g(\gamma)|^2=h_g(\gamma_\rho)$, so under RH $\ft_T[g]$ is exactly the
contribution of the window to the zero side of the explicit formula.

\begin{theorem}[Second-moment identity]\label{thm:id}
$\displaystyle V(d,T)=\int_{-2d}^{2d}|S(\xi)|^2\,d\xi\ \ge 0$, and
$V(d,T)=4dN+\sum 2\sin\!\bigl(2d(\gamma-\gamma')\bigr)/(\gamma-\gamma')$,
the sum over ordered pairs of distinct zeros of the window, the summand being
read as $4d$ when $\gamma=\gamma'$; the diagonal is $4dN$. In particular $0\le V(d,T)\le4dN^2$.
\end{theorem}
\begin{proof}
For $\gamma\ne\gamma'$ one has
\[
\int_{-2d}^{2d}e^{i(\gamma-\gamma')\xi}\,d\xi=\frac{2\sin(2d(\gamma-\gamma'))}{\gamma-\gamma'},
\]
and the integral is $4d$ for $\gamma=\gamma'$. Summing
over the finite index set and interchanging with the integral
gives $\int_{-2d}^{2d}|S|^2=\sum_{\gamma,\gamma'}\int_{-2d}^{2d}
e^{i(\gamma-\gamma')\xi}d\xi=V(d,T)$. Nonnegativity is manifest from the
left side, and the upper bound follows from $|2\sin(2du)/u|\le4d$.
\end{proof}

The identity holds for any finite multiset of real numbers; it says nothing
specific about zeros. It makes $V$ a band-limited $L^2$ energy of the zero
exponential sum.

\begin{proposition}[Band-limited Parseval bound]\label{prop:parseval}
For $g$ of type $d$,
\[
\ft_T[g]=\int_{-2d}^{2d}(g\star\tilde g)(\xi)\,S(\xi)\,d\xi,
\]
hence
\[
0\le\ft_T[g]\le\|g\star\tilde g\|_{L^2}\,V(d,T)^{1/2}\le2\sqrt d\,\|g\|^2\,V(d,T)^{1/2}.
\]
Consequently a bound $V(d,T)\le B$ yields $\ft_T[g]\le2\sqrt{dB}\,\|g\|^2$.
\end{proposition}
\begin{proof}
Sum the displayed formula for $|\widehat g(r)|^2$ over $r=\gamma$ in the window
(a finite sum) to get the identity. Cauchy--Schwarz on $[-2d,2d]$ and
Theorem~\ref{thm:id} give the first inequality; the second follows from
$|g\star\tilde g|\le\|g\|^2$ on an interval of length $4d$. (By Plancherel,
$\|g\star\tilde g\|_{L^2}=(2\pi)^{-1/2}\|\widehat g\|_{L^4}^2$.)
\end{proof}

\begin{remark}
Since $\ft_T[g]\ge0$, the proposition is an \emph{upper} bound for the window
contribution. Passing from it to a lower bound for the full form $\ft=\fa-\fp$
would require control of the zeros outside the window (and, without RH, of
off-line zeros), which we do not supply.
\end{remark}

\section{A prime-side mean square and the frequency split}
On the prime side, the analogous second moment is that of the Dirichlet
polynomial $P_d(t)=\sum_{n\le e^{2d}}\Lambda(n)n^{-1/2-it}$ over the window.
Expanding the square and integrating,
\begin{equation}\label{eq:exp}
\frac1H\int_T^{T+H}|P_d(t)|^2dt
=\underbrace{\sum_{n\le e^{2d}}\frac{\Lambda(n)^2}{n}}_{\text{diagonal}}
+\underbrace{\sum_{\substack{m\ne n\\ m,n\le e^{2d}}}\frac{\Lambda(m)\Lambda(n)}{\sqrt{mn}}\,
K_{T,H}\!\bigl(\log\tfrac nm\bigr)}_{\text{off-diagonal}},\qquad
K_{T,H}(u)=e^{iTu}\,\frac{e^{iHu}-1}{iHu}.
\end{equation}
This identity is elementary; we do not establish an identity between it and
$V(d,T)$, and use it only to identify the diagonal (Lemma~\ref{lem:diag}).

For the zeros we use Montgomery's normalization~\cite{Mont}. Let $L=\log T$,
$w(u)=4/(4+u^2)$, and
\[
F(\alpha,T)=\frac{2\pi}{T\log T}\sum_{0<\gamma,\gamma'\le T}T^{i\alpha(\gamma-\gamma')}w(\gamma-\gamma'),
\]
the sum over ordinates of zeros with multiplicity. Consider the weighted second
moment over the full range $0<\gamma\le T$,
\[
V_w(d,T)=\sum_{0<\gamma,\gamma'\le T}w(\gamma-\gamma')\,\frac{2\sin\!\bigl(2d(\gamma-\gamma')\bigr)}{\gamma-\gamma'}
=\int_{-2d}^{2d}\sum_{0<\gamma,\gamma'\le T}w(\gamma-\gamma')e^{i(\gamma-\gamma')\xi}d\xi .
\]
Substituting $\xi=\alpha L$ and putting $A=2d/L$ and
$\widehat\phi=\mathbf 1_{[-A,A]}$ (a fixed bounded even weight once $A$ is
fixed), one gets exactly
\[
V_w(d,T)=\frac{TL^2}{2\pi}\int_{\bbR}\widehat\phi(\alpha)F(\alpha,T)\,d\alpha .
\]
Split at the critical frequency $\alpha=1$ (we take $A\ge1$):
\[
V_w(d,T)=\underbrace{\frac{TL^2}{2\pi}\int_{|\alpha|\le1}F}_{\text{(I) low}}
+\underbrace{\frac{TL^2}{2\pi}\int_{1<|\alpha|\le A}F}_{\text{(II) high}} .
\]
Recall $N(T)=\#\{0<\gamma\le T\}\sim TL/2\pi$, so $TL^2/2\pi\sim N(T)L$. We do
not prove a comparison between $V_w$ and the window moment $V$.

\begin{proposition}[Low frequencies, under RH]\label{prop:low}
Assume RH. Then $\mathrm{(I)}=(2+o(1))\,N(T)\log T$; in particular
$\mathrm{(I)}\ll dN(T)$ for $d\ge L/2$.
\end{proposition}
\begin{proof}
By Montgomery's theorem~\cite{Mont} (in the form uniform on $0\le\alpha\le1$,
see~\cite{GM}), assuming RH, $F(\alpha,T)=(1+o(1))T^{-2|\alpha|}\log T+|\alpha|+o(1)$
uniformly for $|\alpha|\le1$ ($F$ is even). Hence
$\int_{-1}^{1}F\,d\alpha=(1+o(1))(1-T^{-2})+1+o(1)=2+o(1)$, and
$\mathrm{(I)}=(2+o(1))TL^2/2\pi=(2+o(1))N(T)L$. Since $4dN(T)\ge2LN(T)$ for
$d\ge L/2$, the second claim follows.
\end{proof}

\begin{proposition}[Conditional window bound]\label{prop:S}
Unconditionally, $0\le V(d,T)\le4dN^2$ (Theorem~\ref{thm:id}). Assume
\emph{Hypothesis $H_c$}: $V(d,T)\ll dN(\log T)^c$ uniformly for windows with
$H\ge1$ and $d\ge\tfrac12\log T$, with a fixed $c\ge0$. Then by
Proposition~\ref{prop:parseval}, for $g$ of type $d$,
$\ft_T[g]\ll d\,N^{1/2}(\log T)^{c/2}\|g\|^2$.
\end{proposition}
\begin{proof}
Substitute $B=C\,dN(\log T)^c$ into Proposition~\ref{prop:parseval}.
\end{proof}

\begin{remark}
The range of $d$ in $H_c$ cannot be extended to $d=o(\log T)$ with $c=0$:
for $|\xi|\le1/H$ every phase $(\gamma-T-\tfrac H2)\xi$ has modulus at most
$\tfrac12$, so $|S(\xi)|\ge N\cos\tfrac12$, and hence, when $2d\ge1/H$,
$V(d,T)\ge 2N^2\cos^2(\tfrac12)/H\gg N\log T$ for $H\ge H_0$, where $H_0$ is an
absolute constant such that $N\gg H\log T$ for $H\ge H_0$ and $T$ large
(Riemann--von Mangoldt formula). We do not
prove $H_c$ for any $c$. Even with it, Proposition~\ref{prop:S}
controls only the window contribution, and it does \emph{not} give $(\kappa)$:
a form bounded below may have infinitely many negative squares. In~\cite{RP16}
the number of near-critical directions of the prime/archimedean pencil is
observed numerically, for $d\le6$, to grow with $d$.
\end{remark}

\section{The high band under RH}
\begin{proposition}[High band, under RH]\label{thm:ae}
Assume RH. Then $\int_c^{c+1}F(\alpha,T)\,d\alpha\le\tfrac83+o(1)$ uniformly in
$c\in\bbR$. Consequently, for every fixed $A\ge1$ and $d=AL/2$,
$\mathrm{(II)}\ll AN(T)L$ and $V_w(d,T)\ll dN(T)$, with an implied constant
independent of $A$.
\end{proposition}
\begin{proof}
This is Goldston's argument~\cite{Gold}. Put $K(\alpha)=\max(0,1-|\alpha|)$,
whose Fourier transform $\int K(\alpha)e^{iu\alpha}d\alpha=(\sin(u/2)/(u/2))^2$
is nonnegative. Then
\[
\int_\bbR F(\alpha,T)K(\alpha-c)\,d\alpha=\frac{2\pi}{TL}\sum_{0<\gamma,\gamma'\le T}
w(\gamma-\gamma')\,e^{icL(\gamma-\gamma')}\Bigl(\frac{\sin(u/2)}{u/2}\Bigr)^2,
\qquad u=L(\gamma-\gamma'),
\]
and since $w\ge0$ the right side is at most its value at $c=0$. By Montgomery's
theorem (as in the proof of Proposition~\ref{prop:low}) that value is
$\int_{-1}^1F(\alpha)(1-|\alpha|)\,d\alpha=1+2\int_0^1\alpha(1-\alpha)\,d\alpha+o(1)
=\tfrac43+o(1)$. Since $F\ge0$ under RH (it is $\frac{2\pi}{TL}\int$ of a squared
modulus by $w(u)=\int e^{-2|v|}e^{iuv}dv$) and $K(\alpha-c-\tfrac12)\ge\tfrac12$ on
$[c,c+1]$, we get $\int_c^{c+1}F\le2(\tfrac43+o(1))$. Covering
$\{1<|\alpha|\le A\}$ by at most $2\lceil A-1\rceil\le2A$ unit intervals gives
$\int_{1<|\alpha|\le A}F\le 6A$ for large $T$, so $\mathrm{(II)}\ll AN(T)L\asymp
dN(T)$; with Proposition~\ref{prop:low} this gives $V_w(d,T)\ll dN(T)$.
\end{proof}

\begin{remark}
The bound needs no hypothesis beyond RH; condition (a) of the next section is
needed only to identify the constant.
\end{remark}

\section{The high band and the variance of primes in short intervals}
\begin{theorem}[Goldston--Montgomery~\cite{GM}]\label{thm:gm}
Assume RH. The following are equivalent:
\begin{enumerate}
\item[(a)] $F(\alpha,T)\sim1$ uniformly for $1\le\alpha\le M$, for every fixed $M$;
\item[(b)] $\displaystyle\int_1^X\bigl(\psi(x+\delta x)-\psi(x)-\delta x\bigr)^2\,dx
\sim\tfrac12\,\delta X^2\log\tfrac1\delta$ uniformly for $X^{-1}\le\delta\le X^{-\varepsilon}$,
for every fixed $\varepsilon>0$.
\end{enumerate}
\end{theorem}

\begin{corollary}[Asymptotic under the uniform form factor]\label{cor:noncirc}
Assume RH and condition \emph{(a)} of Theorem~\ref{thm:gm}, equivalently
condition \emph{(b)}. Then for every fixed $A\ge1$ and $d=AL/2$,
$V_w(d,T)=(4+o(1))\,dN(T)$. Thus, under RH, $V_w\ll dN(T)$
(Proposition~\ref{thm:ae}), and under (a), equivalently (b), $V_w=(4+o(1))dN(T)$.
\end{corollary}
\begin{proof}
Since $F$ is even, (a) gives $\int_{1<|\alpha|\le A}F=2(A-1)+o(1)$, and
Proposition~\ref{prop:low} gives $\int_{|\alpha|\le1}F=2+o(1)$. Hence
$V_w=(TL^2/2\pi)(2A+o(1))=(1+o(1))N(T)L\cdot2A=(4+o(1))dN(T)$.
\end{proof}

\section{The diagonal, and the negative index}\label{sec:kappa}
We record the diagonal asymptotic and recall, from~\cite{RP16}, the reduction
for the stronger property~$(\kappa)$.

\begin{lemma}[Diagonal asymptotic]\label{lem:diag}
The diagonal of \eqref{eq:exp} satisfies, unconditionally,
\[
\sum_{n\le e^{2d}}\frac{\Lambda(n)^2}{n}\sim\tfrac12(2d)^2=2d^2\qquad(d\to\infty).
\]
\end{lemma}
\begin{proof}
On primes, $\Lambda(p)^2=(\log p)^2$. The prime powers with exponent at
least $2$ contribute a bounded amount, since
\[
\sum_p\sum_{k\ge2}\frac{(\log p)^2}{p^k}<\infty .
\]
By partial
summation and the prime number theorem $\vartheta(u):=\sum_{p\le u}\log p\sim u$,
we get
\[
\sum_{p\le X}\frac{(\log p)^2}{p}=\int_{2^-}^X\frac{\log u}{u}\,d\vartheta(u)
\sim\int_2^X\frac{\log u}{u}\,du\sim\tfrac12(\log X)^2 .
\]
With $X=e^{2d}$ this is
$2d^2$.
\end{proof}

The negative index of $\ft$ is
\[
\kappa=\#\{\text{negative squares of }\ft\}=\sup_d\kappa_d\in\{0,1,\dots,\infty\},
\]
$\kappa_d$ being the negative index on $PW_d$; RH gives $\kappa=0$. Heuristically
one expects $\kappa$ to be finite exactly when finitely many zeros lie off the
line, and comparable with their number; this is not proved (see~\cite[\S1]{RP16}),
and we use $\kappa$ only as defined.

\begin{proposition}[Relative-form-bound reduction~{\cite[Prop.~2.1]{RP16}}]\label{prop:rfb}
Split the prime form at a cutoff $X$, $\fp=\fp_{\le X}+\fp_{>X}$; the head is a
bounded form, $|\fp_{\le X}[g]|\le B_X\|g\|^2$ with $B_X=2\sum_{n\le X}\Lambda(n)n^{-1/2}$
(it is not of finite rank). Assume $\fa$ is closed and bounded below on a closed
subspace $\mathcal H$ of test functions, with an associated operator $A$ of
compact resolvent. If for some $X$
\[
\fp_{>X}[g]\le a\,\fa[g]+C\|g\|^2\quad(g\in\operatorname{dom}\fa),\qquad 0\le a<1
\qquad\text{(RFB$_X$)},
\]
then $\ft[g]\ge(1-a)\fa[g]-(C+B_X)\|g\|^2$, and $\kappa$ on $\mathcal H$ is at most
the number of eigenvalues $\mu$ of $A$ with $(1-a)\mu<C+B_X$, which is finite.
\end{proposition}

On each $PW_d$ the hypotheses hold with $X=e^{2d}$, so $\kappa_d<\infty$ for
every $d$~\cite[\S2]{RP16}; the open question is uniformity in $d$. With
$\lambda_{\max}(X,d)=\sup_g\fp_{>X}[g]/\fa[g]$ on a finite-dimensional band-$d$
space where $\fa\succ0$, the tail produces no negative directions there once
$\lambda_{\max}(X,d)<1$. Computations of the kind described in~\cite{RP16}, given in
Appendix~\ref{app:rfb}, observe numerically, for $d=4,5$, that the cutoff
$X^\ast(d)$ at which $\lambda_{\max}$ first drops below $1$ grows with the band
(of order $e^{2d}/10$), which, if it persists, means that (RFB$_X$) with $C=0$
cannot hold uniformly in $d$ for any fixed $X$; and, for $d=3,\dots,6$, that the number
of generalized eigenvalues of $(\fp,\fa)$ near $1$ grows roughly linearly with $d$
(Appendix~\ref{app:near}). These observations do not decide whether $\kappa$ is
finite.

\section{Summary of the hierarchy}
The status of the steps discussed above is as follows.
\begin{itemize}
\item Unconditionally: the identity of Theorem~\ref{thm:id}, the trivial bound
$V\le4dN^2$, and the Parseval bound $\ft_T[g]\le2\sqrt d\|g\|^2V^{1/2}$ for the
window contribution (Proposition~\ref{prop:parseval}).
\item Under Hypothesis $H_c$: $\ft_T[g]\ll dN^{1/2}(\log T)^{c/2}\|g\|^2$
(Proposition~\ref{prop:S}).
\item Under RH: the low band is $(2+o(1))N(T)\log T$ (Proposition~\ref{prop:low});
the high band is $O(dN(T))$, so $V_w\ll dN(T)$ (Proposition~\ref{thm:ae}); and the uniform high-band asymptotic, equivalent by
Goldston--Montgomery to the short-interval prime variance asymptotic, gives
$V_w\sim4dN$ (Corollary~\ref{cor:noncirc}).
\item From these window bounds to (S) for the full form, from (S) to ($\kappa$)
(for which~\cite{RP16} gives a conditional reduction), and from ($\kappa$) to the
sign of $\ft$, nothing is proved here.
\end{itemize}

\section{Conclusion}
The second moment $V(d,T)$ of the zero exponential sum is an exact band-limited
$L^2$ energy (Theorem~\ref{thm:id}) and bounds the window contribution to the
Weil form (Proposition~\ref{prop:parseval}). Assuming RH, Montgomery's theorem
gives the low-frequency part of its weighted version; the high-frequency part
is $O(dN(T))$ under RH alone (Proposition~\ref{thm:ae}), and its asymptotic,
giving $V_w\sim4dN(T)$, is conditional on the uniform form factor, which by Goldston--Montgomery is equivalent,
under RH, to an asymptotic for the variance of primes in short intervals
(Theorem~\ref{thm:gm}, Corollary~\ref{cor:noncirc}). We make no claim about the
sign of $\ft$, and hence none about RH.

\appendix
\section{Numerical check of the second-moment identity}\label{app:code}
The following program computes both sides of $V(d,T)=\int_{-2d}^{2d}|S|^2$ for
the ordinates of the zeros $\#1000$--$\#1079$ of $\zeta$ and the ratio
$V/(dN)$. Since Theorem~\ref{thm:id} holds for any real numbers, the comparison
measures only the error of the Riemann-sum quadrature.

\lstinputlisting[language=Python,caption={\texttt{code.py}}]{code.py}

\medskip\noindent\textbf{Output} (stored in \texttt{results\_code.txt}; $N=80$,
$T\approx1465$, $\alpha=2d/L$ with $L=\log T$, as in Section~3):
\begin{lstlisting}
      d   alpha=2d/L    V (pairsum)   V (integral)   rel.err    V/(dN)
   0.25       0.0686       429.5252       429.5253   9.4e-08   21.4763
   0.50       0.1372       442.1774       442.1776   3.0e-07   11.0544
   0.75       0.2058       457.2294       457.2296   4.0e-07    7.6205
   1.00       0.2744       484.6011       484.6012   3.3e-07    6.0575
   1.50       0.4115       552.7181       552.7184   4.8e-07    4.6060
   2.00       0.5487       624.1183       624.1195   1.9e-06    3.9007
   3.00       0.8231       911.2829       911.2857   3.1e-06    3.7970
\end{lstlisting}
The two sides of Theorem~\ref{thm:id} agree to within $3.1\times10^{-6}$. For
this single window $V/(dN)$ decreases to about $3.8$ at $d=3$; this is
consistent with, but no evidence for, a uniform bound.

\section{The tail relative form bound \texorpdfstring{$\lambda_{\max}(X,d)$}{lambdamax}}\label{app:rfb}
The following program computes the extreme generalized eigenvalues of the pencil
$(\fp_{>X},\fa)$ on a band-$d$ space spanned by $21$ functions whose Fourier
transforms are sinc functions at spacing $\pi/d$ centred near $T_0=10^3$, with
the archimedean form taken without its pole terms, restricting the prime form to
the tail $n>X$.
\lstinputlisting[language=Python,caption={\texttt{atail.py}}]{atail.py}
\noindent\textbf{Output} (stored in \texttt{results\_atail.txt}). For $d=4,5$,
$\lambda_{\max}(1,d)=1.00000$ to the printed precision. On the grid of cutoffs
tested, $\lambda_{\max}(X,d)$ first drops below $1$ between $X=150$ and $X=400$
for $d=4$, and between $X=1000$ and $X=5506$ for $d=5$; this is consistent with a
cutoff $X^\ast(d)$ of order $e^{2d}/10$, a numerical observation on two values
of $d$ only.

\section{Near-critical directions of the full pencil}\label{app:near}
The following program computes the spectrum of the full pencil $(\fp,\fa)$ on
band-$d$ spaces of dimension $6d+13$ at $T_0=10^3$, for $d=3,\dots,6$, with the archimedean form taken without its pole terms, as in Appendix~\ref{app:rfb}.
\lstinputlisting[language=Python,caption={\texttt{near.py}}]{near.py}
\noindent\textbf{Output} (stored in \texttt{results\_near.txt}):
$\#\{\lambda>0.99\}=3,11,19,27$ for $d=3,4,5,6$, the gap
$\lambda_{\max}-\lambda_2$ is at most $1.3\times10^{-4}$, and the fraction of
eigenvalues above $0.99$ rises from about $10\%$ to $55\%$. This is a numerical
observation for four values of $d$; it does not determine the behaviour as
$d\to\infty$.

\begin{thebibliography}{9}
\bibitem{RP16} D.~A.~Trejo Pizzo, \emph{The off-line zeros as a Pontryagin index: a conditional reduction and numerical observations on the prime tail}, preprint (Paper 11).
\bibitem{GM} D.~A.~Goldston and H.~L.~Montgomery, \emph{Pair correlation of zeros and primes in short intervals}, in \emph{Analytic Number Theory and Diophantine Problems} (Stillwater, 1984), Progr.\ Math.\ \textbf{70}, Birkh\"auser, Boston, 1987, pp.~183--203.
\bibitem{Mont} H.~L.~Montgomery, \emph{The pair correlation of zeros of the zeta function}, in \emph{Analytic Number Theory}, Proc.\ Sympos.\ Pure Math.\ \textbf{24}, Amer.\ Math.\ Soc., Providence, 1973, pp.~181--193.
\bibitem{Gold} D.~A.~Goldston, \emph{On the pair correlation conjecture for zeros of the Riemann zeta-function}, J.\ reine angew.\ Math.\ \textbf{385} (1988), 24--40.
\bibitem{Weil} A.~Weil, \emph{Sur les ``formules explicites'' de la th\'eorie des nombres premiers}, Comm.\ S\'em.\ Math.\ Univ.\ Lund (1952), tome suppl\'ementaire, 252--265.
\end{thebibliography}

\end{document}
